% =====================================================================
% CHƯƠNG 3 - NGHIÊN CỨU VÀ TỐI ƯU MÔ HÌNH YOLO
% =====================================================================
\chapter{NGHIÊN CỨU VÀ TỐI ƯU MÔ HÌNH YOLO}
\label{chap:toi-uu}

Chương này trình bày phần nghiên cứu mô hình trước khi đưa model vào nền tảng.
Đồ án sử dụng YOLO26n --- biến thể nhỏ nhất của thế hệ YOLO26 --- làm điểm xuất
phát và xây dựng một pipeline tối ưu nối tiếp gồm sáu giai đoạn: tìm siêu tham số
huấn luyện baseline bằng thuật toán Sea-Horse Optimizer (SHO), huấn luyện thưa
hóa BatchNorm, cắt tỉa kênh có cấu trúc, phục hồi độ chính xác bằng chưng cất tri
thức hoặc fine-tune, huấn luyện nhận biết lượng tử hóa (QAT) và chuyển đổi sang
TensorRT INT8. Mỗi giai đoạn nhận checkpoint của giai đoạn trước và sinh
checkpoint cho giai đoạn sau.

\section{Bài toán tối ưu}

\subsection{Bài toán ứng dụng}

Bài toán ứng dụng là giám sát tuân thủ trang bị bảo hộ lao động (PPE --
Personal Protective Equipment) trên công trường xây dựng: phát hiện người lao động
và các trang bị mũ bảo hộ, áo phản quang, găng tay, ủng, kính bảo hộ, đồng thời
phát hiện trường hợp \emph{thiếu} trang bị. Camera công trường thường ở xa hạ tầng
mạng tốt, nên suy luận tại biên trên thiết bị Jetson là lựa chọn phù hợp.

\subsection{Mô hình baseline}

Mô hình baseline là YOLO26n do Ultralytics phát hành \cite{ultralytics}, khởi tạo từ trọng số tiền
huấn luyện trên COCO rồi huấn luyện trên tập PPE ở độ phân giải $640\times640$,
suy luận FP32. YOLO26n có 2{,}57 triệu tham số và 6{,}12 GFLOPs tại $640\times640$
(đo trên checkpoint gốc). Baseline cung cấp mốc chất lượng, kích thước và hiệu
năng để đánh giá mọi biến thể sau đó.

\subsection{Hạn chế của mô hình baseline}

Dù đã là biến thể nhỏ, baseline chưa được tối ưu riêng cho Jetson: mọi kênh
tích chập đều được giữ lại kể cả kênh đóng góp rất ít, và phép tính chạy ở FP32
trong khi GPU Jetson có đơn vị tính FP16/INT8 nhanh hơn nhiều. Trên Jetson Nano
(GPU Maxwell 128 nhân, không có Tensor Core), chi phí FP32 là nút thắt chính; trên
Jetson Orin Nano, Tensor Core thế hệ Ampere chỉ phát huy khi mô hình chạy INT8
hoặc FP16. Mức hạn chế cụ thể được xác nhận bằng số đo trên thiết bị thay vì chỉ
suy luận từ FLOPs.

\subsection{Mục tiêu tối ưu}

Đồ án tối ưu đồng thời sáu tiêu chí: độ chính xác, latency, FPS, kích thước mô
hình, bộ nhớ GPU và công suất. Bài toán được biểu diễn dưới dạng:
\begin{equation}
  \min_{m\in\mathcal{M}}
  \bigl(-A(m),\ L(m),\ -F(m),\ S(m),\ G(m),\ P(m)\bigr),
\end{equation}
trong đó $m$ là một cấu hình model, $A$ là $mAP@50{:}95$, $L$ là latency, $F$ là
FPS, $S$ là kích thước, $G$ là bộ nhớ GPU và $P$ là công suất. Thay vì gộp tùy ý
thành một điểm duy nhất, đồ án loại trước các cấu hình vi phạm ràng buộc rồi so
sánh các cấu hình còn lại theo trade-off.

Một cấu hình chỉ được xem là ứng viên triển khai khi thỏa đồng thời ba ràng buộc:
\begin{enumerate}
  \item $mAP@50{:}95$ trên tập test giảm không quá \textbf{2 điểm phần trăm} so
        với baseline FP32;
  \item latency suy luận p95 không vượt \textbf{66~ms}, tương ứng FPS mục tiêu
        15 --- giá trị mặc định \texttt{target\_fps} của AI Runtime;
  \item thiết bị chạy ổn định trong bài đo kéo dài 30 phút, không giảm xung do
        nhiệt làm FPS rơi dưới ngưỡng.
\end{enumerate}

\section{Xây dựng mô hình YOLO baseline}

\subsection{Dataset}

Tập dữ liệu sử dụng là \textbf{Construction-PPE} do Ultralytics công bố (khoảng
178~MB), gồm ảnh công trường với nhãn hộp bao cho 11 lớp. Tập được chia sẵn thành
train/validation/test như Bảng~\ref{tab:ppe-dataset}, tỷ lệ xấp xỉ 80/10/10.

\begin{table}[H]
\centering
\caption{Tập dữ liệu Construction-PPE}
\label{tab:ppe-dataset}
\begin{tabularx}{\textwidth}{L{3.2cm}Y}
\toprule
\textbf{Thuộc tính} & \textbf{Giá trị}\\
\midrule
Số ảnh & 1.416 (train 1.132, validation 143, test 141)\\
Số lớp & 11\\
Lớp trang bị & \texttt{helmet}, \texttt{gloves}, \texttt{vest}, \texttt{boots}, \texttt{goggles}\\
Lớp thiếu trang bị & \texttt{no\_helmet}, \texttt{no\_goggle}, \texttt{no\_gloves}, \texttt{no\_boots}\\
Lớp khác & \texttt{Person}, \texttt{none}\\
Định dạng nhãn & YOLO (lớp, tâm $x$, tâm $y$, rộng, cao đã chuẩn hóa)\\
\bottomrule
\end{tabularx}
\end{table}

Ngoài tập công khai, đồ án thu thập thêm dữ liệu thực tế gồm các đoạn video và
ảnh công trường (thư mục \texttt{dataset\_raw}) và một tập PPE tùy chỉnh 9 lớp
xuất từ Roboflow (tập validation 229 ảnh) để đánh giá khả năng tổng quát trên dữ
liệu ngoài phân phối huấn luyện. Nhãn được kiểm tra về chỉ số lớp, hộp vượt biên
và ảnh lỗi trước khi huấn luyện; các khung trích từ cùng một video được giữ trong
cùng một tập để tránh rò rỉ ngữ cảnh giữa train và test.

\subsection{Tiền xử lý và tăng cường dữ liệu}

Ảnh được đưa về $640\times640$ bằng letterbox (giữ tỷ lệ, đệm viền), chuyển
BGR$\rightarrow$RGB và chuẩn hóa về $[0,1]$. Các phép tăng cường dùng bộ tăng
cường của Ultralytics: biến đổi HSV, xoay, tịnh tiến, co giãn, cắt xiên, phối
cảnh, lật dọc/ngang, mosaic, mixup và copy-paste. Thay vì dùng hệ số mặc định,
cường độ của các phép này được tìm tự động bằng SHO cùng các siêu tham số tối ưu
hóa (mục tiếp theo). Mosaic được tắt ở 10 epoch cuối (\texttt{close\_mosaic=10})
để mô hình thích nghi với phân bố ảnh thật. Cấu hình tiền xử lý suy luận
(letterbox, thứ tự kênh, chuẩn hóa) được giữ nguyên khi export và trong AI Runtime.

\subsection{Tìm siêu tham số bằng Sea-Horse Optimizer}

Sea-Horse Optimizer \cite{zhao2023sho} là thuật toán tối ưu bầy đàn mô phỏng
hành vi di chuyển, săn mồi và sinh sản của cá ngựa. Mỗi cá thể là một vector
siêu tham số; ở mỗi thế hệ, cá thể di chuyển theo quỹ đạo xoắn ốc hoặc bước
Lévy quanh cá thể tốt nhất (khám phá), săn mồi với xác suất thành công ngẫu
nhiên (khai thác), rồi nửa quần thể tốt nhất sinh thế hệ con. Đồ án hiện thực SHO
độc lập với Ultralytics (\texttt{optimization/sho.py}) để kiểm thử quy tắc cập
nhật bằng hàm mục tiêu toán học trước khi dùng cho YOLO.

Không gian tìm kiếm gồm 18 chiều (Bảng~\ref{tab:sho-space}). Hàm mục tiêu cần
cực tiểu là $f(\mathbf{x}) = 1 - mAP@50{:}95(\mathbf{x})$, trong đó
$mAP@50{:}95(\mathbf{x})$ đo trên tập validation sau một lần huấn luyện \emph{proxy}
ngắn 3 epoch từ cùng checkpoint khởi tạo. Ngân sách tìm kiếm: quần thể 6 cá thể,
3 thế hệ. Bộ siêu tham số tốt nhất được ghi ra \texttt{optimized\_baseline.yaml}
và dùng để huấn luyện baseline đầy đủ.

\begin{table}[H]
\centering
\caption{Không gian tìm kiếm siêu tham số của SHO}
\label{tab:sho-space}
\begin{tabularx}{\textwidth}{L{3.4cm}C{2.8cm}Y}
\toprule
\textbf{Tham số} & \textbf{Khoảng} & \textbf{Ý nghĩa}\\
\midrule
\texttt{lr0} & $[10^{-5}, 10^{-1}]$ & Learning rate ban đầu\\
\texttt{lrf} & $[0{,}01; 1]$ & Tỷ lệ learning rate cuối/ban đầu\\
\texttt{momentum} & $[0{,}6; 0{,}98]$ & Momentum/$\beta_1$\\
\texttt{weight\_decay} & $[0; 0{,}001]$ & Hệ số suy giảm trọng số\\
\texttt{warmup\_epochs} & $[0; 5]$ & Số epoch khởi động\\
\texttt{warmup\_momentum} & $[0; 0{,}95]$ & Momentum giai đoạn khởi động\\
\texttt{hsv\_h}, \texttt{hsv\_s}, \texttt{hsv\_v} & $[0; 0{,}9]$ & Biến đổi màu, bão hòa, độ sáng\\
\texttt{degrees} & $[0; 45]$ & Góc xoay (độ)\\
\texttt{translate}, \texttt{scale} & $[0; 0{,}9]$ & Tịnh tiến, co giãn\\
\texttt{shear} & $[0; 10]$ & Cắt xiên (độ)\\
\texttt{perspective} & $[0; 0{,}001]$ & Phối cảnh\\
\texttt{flipud} & $[0; 1]$ & Xác suất lật dọc\\
\texttt{mosaic}, \texttt{mixup}, \texttt{copy\_paste} & $[0; 1]$ & Xác suất các phép ghép ảnh\\
\bottomrule
\end{tabularx}
\end{table}

\subsection{Huấn luyện}

Baseline được huấn luyện 100 epoch, batch 16, ảnh $640\times640$, trên GPU
Google Colab; toàn bộ đầu ra được ghi trực tiếp lên Google Drive để không mất khi
phiên Colab bị ngắt. Optimizer và learning rate lấy từ kết quả SHO. Seed được cố
định và \texttt{deterministic=True}. Checkpoint tốt nhất được chọn theo
$mAP@50{:}95$ trên tập validation, không chọn theo tập test.

\subsection{Đánh giá baseline}

\begin{table}[H]
\centering
\caption{Kết quả chất lượng của mô hình baseline trên tập test Construction-PPE}
\label{tab:baseline-quality}
\resizebox{\textwidth}{!}{%
\begin{tabular}{lcccccc}
\toprule
\textbf{Model} & \textbf{Params} & \textbf{GFLOPs} & \textbf{Precision} & \textbf{Recall} & \textbf{mAP@50} &
\textbf{mAP@50:95}\\
\midrule
YOLO26n baseline FP32 & 2,57 M & 6,12 & \cbs{} & \cbs{} & \cbs{} & \cbs{}\\
\bottomrule
\end{tabular}}
\end{table}

Ngoài chỉ số tổng hợp, AP theo lớp và confusion matrix được dùng để nhận diện
lớp yếu. Các lớp ``thiếu trang bị'' (\texttt{no\_gloves}, \texttt{no\_goggle}) có
đối tượng nhỏ và ít mẫu, dự kiến là nhóm nhạy cảm nhất với cắt tỉa và lượng tử
hóa, nên được theo dõi riêng ở các giai đoạn sau.

\section{Tối ưu bằng Pruning}

\subsection{Lựa chọn phương pháp pruning}

Đồ án chọn \textbf{cắt tỉa kênh có cấu trúc} theo phương pháp Network Slimming
\cite{liu2017slimming}, vì kết quả là một mạng nhỏ hơn thật sự (ít kênh hơn) mà
TensorRT khai thác trực tiếp, không cần phần cứng hỗ trợ ma trận thưa. Ý tưởng
là dùng hệ số co giãn $\gamma$ của lớp BatchNorm sau mỗi tích chập làm thước đo
tầm quan trọng của kênh: kênh có $|\gamma|$ nhỏ đóng góp ít vào đầu ra và có thể
loại bỏ. Unstructured pruning chỉ làm thưa trọng số mà không giảm kích thước
tensor nên không đem lại tăng tốc trên Jetson, và không được chọn.

\subsection{Huấn luyện thưa hóa}

Để phân bố $\gamma$ tách rõ giữa kênh quan trọng và kênh thừa, mô hình được
huấn luyện với phạt L1 trên $\gamma$:
\begin{equation}
  \mathcal{L} = \mathcal{L}_{det} + s \sum_{\gamma \in \Gamma} |\gamma|,
\end{equation}
hiện thực bằng cách cộng $s\cdot\mathrm{sign}(\gamma)$ vào gradient của $\gamma$
sau mỗi lần lan truyền ngược. Hệ số $s$ giảm dần theo epoch,
$s_e = s\,(1 - 0{,}9\,e/E)$, để phạt mạnh ở đầu và để mô hình phục hồi độ chính
xác ở cuối. Cấu hình: $s = 10^{-2}$, 200 epoch, ảnh 640. Một callback in thống kê
$\gamma$ mỗi epoch để theo dõi mức thưa. Thử nghiệm với $s$ hằng số cho thấy độ
chính xác validation sụp gần về 0, nên lịch giảm dần là cần thiết.

\subsection{Cắt tỉa theo ngưỡng toàn cục}

Script \texttt{prune.py} thực hiện:
\begin{enumerate}
  \item Thu thập mọi lớp BN, trừ các lớp trong danh sách bảo vệ (nhánh có ràng
        buộc cộng dư, lớp đầu vào của head và các tầng nhạy cảm);
  \item Sắp xếp toàn bộ $|\gamma|$ và lấy ngưỡng $\tau$ tại vị trí
        \texttt{prune\_ratio}. Ngưỡng bị chặn trên bởi giá trị nhỏ nhất trong các
        $\max|\gamma|$ của từng lớp để không lớp nào bị cắt hết kênh; nếu tỷ lệ yêu
        cầu vượt giới hạn này, tỷ lệ tự động được hạ xuống;
  \item Tạo mặt nạ kênh $|\gamma| \ge \tau$ cho từng lớp và dựng kiến trúc mới
        với các module chuyên dụng (\texttt{C3k2Pruned}, \texttt{C2PSAPruned},
        \texttt{SPPFPruned}, \texttt{DetectPruned}) theo cấu hình
        \texttt{yolo26-pruned.yaml};
  \item Sao chép trọng số của các kênh được giữ sang mô hình mới, lưu mô hình
        (lớp \texttt{DetectionModelPruned}) cùng từ điển mặt nạ kênh.
\end{enumerate}

Tỷ lệ cắt tỉa chính được chọn là \textbf{30\%}. Với YOLO26n, cắt 30\% kênh BN
cho kết quả ở Bảng~\ref{tab:pruning-complexity}: số tham số giảm 15,7\% và
GFLOPs giảm 21,2\%. Mức giảm FLOPs lớn hơn mức giảm tham số vì các kênh bị cắt
tập trung ở tầng có độ phân giải không gian lớn.

\begin{table}[H]
\centering
\caption{Độ phức tạp của mô hình trước và sau cắt tỉa 30\% (đo tại $640\times640$)}
\label{tab:pruning-complexity}
\begin{tabular}{lccc}
\toprule
\textbf{Mô hình} & \textbf{Tham số} & \textbf{GFLOPs} & \textbf{Thay đổi}\\
\midrule
YOLO26n baseline & 2.572.280 & 6,12 & --\\
YOLO26n pruned 30\% & 2.169.072 & 4,82 & $-15{,}7\%$ tham số, $-21{,}2\%$ FLOPs\\
\bottomrule
\end{tabular}
\end{table}

\subsection{Phục hồi độ chính xác}

Ngay sau khi cắt, độ chính xác giảm mạnh và cần được phục hồi. Đồ án so sánh hai
phương án thay thế nhau (không chạy nối tiếp):

\textbf{Fine-tune thường:} huấn luyện tiếp mô hình đã cắt 100 epoch, AdamW,
$lr_0 = 10^{-3}$. Cấu trúc đã cắt được khóa bởi kiến trúc mới nên kênh bị loại
không xuất hiện trở lại.

\textbf{Chưng cất tri thức:} mô hình đã cắt là \emph{student}, YOLO26l (hoặc
YOLO26x) là \emph{teacher}. Hàm mất mát
\begin{equation}
  \mathcal{L} = \mathcal{L}_{det} + \lambda_{cwd}\,\mathcal{L}_{CWD}
  + \alpha\,(\mathcal{L}_{cls\text{-}KD} + \mathcal{L}_{reg\text{-}KD}),
\end{equation}
trong đó $\mathcal{L}_{CWD}$ là Channel-Wise Distillation \cite{shu2021cwd}: với
mỗi kênh của bản đồ đặc trưng, phân bố không gian được chuẩn hóa bằng softmax với
nhiệt độ $T$ và student được ép khớp phân bố của teacher bằng KL divergence; một
lớp tích chập $1\times1$ (adapter) đưa số kênh của student về bằng teacher.
$\mathcal{L}_{cls\text{-}KD}$ chưng cất điểm phân lớp ở head. Cấu hình mặc định:
$\lambda_{cwd}=50$, $T=4$; trọng số $\alpha$ tăng dần theo epoch. Vì hàm mất mát
đã bao gồm $\mathcal{L}_{det}$ trên nhãn thật, chưng cất thay thế hoàn toàn bước
fine-tune.

Kết quả thực nghiệm của các phương án được trình bày ở Mục~\ref{sec:pilot-results}.

\section{Tối ưu bằng QAT}

\subsection{Thiết lập quantization}

QAT được hiện thực bằng thư viện \texttt{pytorch-quantization} của NVIDIA, là
đường chính thức để tạo mô hình INT8 có thông tin lượng tử tường minh cho
TensorRT \cite{jacob2018}. Hàm \texttt{quant\_modules.initialize()} thay mọi
\texttt{nn.Conv2d} bằng \texttt{QuantConv2d} trước khi dựng mô hình. Theo mặc định
của thư viện, lượng tử hóa là \textbf{đối xứng}: trọng số lượng tử \textbf{theo
từng kênh đầu ra} (per-channel), activation lượng tử \textbf{theo tensor}
(per-tensor). Với các phép cộng dư và ghép kênh trong các khối
\texttt{C3k2}, \texttt{C2PSA} đã cắt tỉa, đồ án viết các lớp lượng tử riêng
(\texttt{quant\_ops\_pruned\_yolo26.py}) chèn \texttt{TensorQuantizer} cho từng
đầu vào, để TensorRT có thể hợp nhất phép toán ở INT8 thay vì chuyển đổi qua lại
giữa INT8 và FP16. Các lượng tử của những nhánh không cần thiết được vô hiệu hóa
để giữ độ chính xác.

\subsection{Hiệu chỉnh}

Trước khi huấn luyện, thang lượng tử (\texttt{amax}) của activation được hiệu
chỉnh trên 256 batch dữ liệu huấn luyện bằng phương pháp \textbf{entropy}: với
mỗi tensor, thu histogram giá trị và chọn ngưỡng cắt cực tiểu hóa KL divergence
giữa phân bố gốc và phân bố sau lượng tử. Phương pháp này bỏ qua các giá trị
ngoại lai hiếm gặp, cho độ phân giải lượng tử tốt hơn so với lấy giá trị lớn nhất.

\subsection{Fake quantization và huấn luyện}

Sau hiệu chỉnh, các nút fake-quantization mô phỏng phép làm tròn và cắt về INT8
trong lan truyền thuận, còn gradient đi qua bằng Straight-Through Estimator. Mô
hình được fine-tune nhẹ 20 epoch, batch 8, optimizer SGD momentum 0,9; learning
rate đầu vào $10^{-3}$ được chia 100 bên trong trainer (thành $10^{-5}$) vì QAT
chỉ cần điều chỉnh nhỏ quanh điểm đã hội tụ. Tùy chọn hiệu chỉnh lại định kỳ
(\texttt{recalib\_every}) mặc định tắt. QAT chỉ được chạy trên checkpoint \emph{đã
phục hồi} sau cắt tỉa, không chạy trên checkpoint vừa cắt còn yếu.

\subsection{Export INT8}

Checkpoint QAT được export sang ONNX, trong đó mỗi tensor lượng tử được biểu diễn
bằng một cặp nút Q/DQ (\texttt{QuantizeLinear} và \texttt{DequantizeLinear})
mang thang lượng tử đã học. TensorRT đọc các nút Q/DQ này để build engine INT8 trực tiếp, không cần bước
hiệu chỉnh PTQ riêng. Thư viện \texttt{pytorch-quantization} yêu cầu CUDA nên giai
đoạn này chạy trên môi trường Linux có GPU (Google Colab), không chạy được trên
máy Windows chỉ có CPU.

\subsection{Đánh giá mô hình QAT}

Mô hình QAT được đánh giá ở ba điểm: checkpoint fake-quant trong PyTorch, ONNX
sau export và TensorRT engine trên thiết bị. Chênh lệch giữa các điểm cho biết sai
số đến từ lượng tử, từ export hay từ backend.

\section{Pipeline tối ưu kết hợp}

\subsection{Quy trình}

Quy trình kết hợp được thể hiện trong Hình~\ref{fig:pruning-qat-pipeline} và
được đóng gói thành notebook \texttt{yolo26\_ppe\_train\_pipeline.ipynb} chạy trên
Google Colab: mỗi ô thực hiện một giai đoạn bằng một script độc lập trong
\texttt{scripts/}, checkpoint của giai đoạn trước là đầu vào của giai đoạn sau.

\begin{figure}[H]
\centering
\resizebox{\textwidth}{!}{%
\begin{tikzpicture}[node distance=0.5cm]
  \node[flowbox] (sho) {SHO search\\+ baseline\\\scriptsize sho\_search.py};
  \node[flowbox, right=of sho] (sp) {Sparsity\\training\\\scriptsize train\_sparsity.py};
  \node[flowbox, right=of sp] (pr) {Channel\\pruning 30\%\\\scriptsize prune.py};
  \node[flowbox, right=of pr] (rec) {Distill hoặc\\fine-tune\\\scriptsize distill.py};
  \node[flowbox, right=of rec] (qat) {QAT\\INT8\\\scriptsize qat.py};
  \node[flowbox, right=of qat, fill=green!6] (trt) {ONNX Q/DQ\\$\rightarrow$ TensorRT};
  \draw[flowarrow] (sho) -- (sp);
  \draw[flowarrow] (sp) -- (pr);
  \draw[flowarrow] (pr) -- (rec);
  \draw[flowarrow] (rec) -- (qat);
  \draw[flowarrow] (qat) -- (trt);
\end{tikzpicture}}
\caption{Pipeline tối ưu YOLO26n từ baseline đến TensorRT INT8}
\label{fig:pruning-qat-pipeline}
\end{figure}

Thứ tự các giai đoạn có chủ đích: cắt tỉa thay đổi cấu trúc nên phải làm trước;
phục hồi độ chính xác trước khi lượng tử để QAT bắt đầu từ một mô hình tốt;
QAT là bước cuối vì mọi thay đổi trọng số sau QAT sẽ làm thang lượng tử mất
hiệu lực. Bảng~\ref{tab:pipeline-hparams} tóm tắt cấu hình từng giai đoạn.

\begin{table}[H]
\centering
\caption{Cấu hình các giai đoạn trong pipeline tối ưu}
\label{tab:pipeline-hparams}
\begin{tabularx}{\textwidth}{L{3.2cm}Y}
\toprule
\textbf{Giai đoạn} & \textbf{Cấu hình}\\
\midrule
SHO + baseline & Quần thể 6, 3 thế hệ, proxy 3 epoch; huấn luyện đầy đủ 100 epoch, batch 16, ảnh 640\\
Sparsity & 200 epoch, $s=10^{-2}$ giảm dần tuyến tính về $0{,}1s$\\
Pruning & Tỷ lệ 0,30, ngưỡng toàn cục trên $|\gamma|$, mô hình \texttt{n}\\
Phục hồi & Chưng cất: teacher YOLO26l, $\lambda_{cwd}=50$, $T=4$, 100 epoch; hoặc fine-tune 100 epoch\\
QAT & 20 epoch, batch 8, SGD, $lr_0=10^{-3}/100$, hiệu chỉnh entropy 256 batch\\
Export & ONNX Q/DQ, ảnh $640\times640$, batch 1; TensorRT FP16 và INT8\\
\bottomrule
\end{tabularx}
\end{table}

\subsection{Ma trận cấu hình đánh giá}

\begin{table}[H]
\centering
\caption{Ma trận cấu hình tối ưu mô hình}
\label{tab:optimization-configs}
\begin{tabular}{lccc}
\toprule
\textbf{Cấu hình} & \textbf{Pruning} & \textbf{Phục hồi / QAT} & \textbf{Precision đích}\\
\midrule
Baseline & 0\% & -- & FP32, FP16\\
Pruned-30 + fine-tune & 30\% & Fine-tune & FP16\\
Pruned-30 + distill & 30\% & Chưng cất & FP16\\
Baseline + QAT & 0\% & QAT & INT8\\
Pruned-30 + distill + QAT & 30\% & Chưng cất + QAT & INT8\\
\bottomrule
\end{tabular}
\end{table}

Mỗi cấu hình dùng cùng dataset split, cùng kích thước ảnh và cùng quy trình
đánh giá. Tên run, seed, checkpoint nguồn và cấu hình được lưu (tệp
\texttt{args.yaml}, \texttt{results.csv} của mỗi run) để truy ngược artifact về
thí nghiệm tạo ra nó.

\section{Chuyển đổi sang TensorRT}

\subsection{Export ONNX}

Mỗi checkpoint ứng viên được export sang ONNX với đầu vào tĩnh
$1\times3\times640\times640$ (batch 1, phù hợp xử lý luồng camera từng khung)
bằng chức năng export của Ultralytics. ONNX checker xác nhận đồ thị hợp lệ; suy
luận đối chiếu giữa PyTorch và ONNX Runtime trên một tập ảnh cố định đo sai khác
tensor đầu ra và sai khác detection cuối. YOLO26 có head end-to-end không cần NMS
nên đồ thị ONNX không chứa bước NMS phụ thuộc dữ liệu, giúp chuyển đổi TensorRT
thuận lợi.

\subsection{TensorRT FP16 và INT8}

Engine được build \emph{trên chính thiết bị đích} vì engine TensorRT gắn với
kiến trúc GPU và phiên bản TensorRT: engine tạo trên Jetson Nano (TensorRT 8.2)
không chạy trên Jetson Orin Nano và ngược lại. Engine FP16 được tạo từ baseline và
các mô hình đã cắt tỉa. Engine INT8 được tạo từ ONNX Q/DQ của mô hình QAT; vì
thang lượng tử đã có trong đồ thị nên không cần tập hiệu chỉnh bổ sung. Log build
ghi lại phiên bản TensorRT, workspace và precision thực tế của từng lớp để xác định
tỷ lệ phép toán thực sự chạy INT8.

\subsection{Kiểm tra độ chính xác sau chuyển đổi}

Kiểm tra gồm ba mức: sai khác tensor trên mẫu cố định, sai khác detection theo
IoU và confidence, và chênh lệch $mAP$ trên toàn tập test. Artifact vi phạm ràng
buộc độ chính xác ở Mục~3.1 bị loại dù latency tốt. Engine đạt yêu cầu được đẩy
lên kho model và triển khai qua nền tảng; AI Runtime nạp trực tiếp tệp
\texttt{.engine} bằng Ultralytics.

\section{Lựa chọn mô hình triển khai}

Model triển khai được chọn theo hai bước: loại các cấu hình vi phạm ràng buộc
độ chính xác, latency hoặc ổn định nhiệt; trong nhóm còn lại, ưu tiên cấu hình có
latency thấp nhất trên thiết bị yếu hơn (Jetson Nano) vì đó là giới hạn của cả đội
thiết bị.

Dựa trên phân tích độ phức tạp và kết quả thực nghiệm sơ bộ (Mục~\ref{sec:pilot-results}),
cấu hình được đề xuất là \textbf{YOLO26n cắt tỉa 30\% + phục hồi + QAT, chạy
TensorRT INT8} (Bảng~\ref{tab:selected-model}). Lý do: cắt tỉa giảm 21\% FLOPs,
INT8 giảm tiếp chi phí mỗi phép toán và bộ nhớ trọng số khoảng 4 lần so với
FP32, còn QAT giữ độ chính xác tốt hơn PTQ. Với thiết bị hỗ trợ kém INT8, engine
FP16 của mô hình đã cắt tỉa là phương án dự phòng. Kết luận cuối cùng phụ thuộc
số đo trên thiết bị ở Chương~6.

\begin{table}[H]
\centering
\caption{Cấu hình model đề xuất triển khai}
\label{tab:selected-model}
\begin{tabularx}{\textwidth}{L{4cm}Y}
\toprule
\textbf{Thuộc tính} & \textbf{Giá trị}\\
\midrule
Kiến trúc & YOLO26n, đầu vào $640\times640$, 11 lớp PPE\\
Phương pháp tối ưu & SHO baseline $\rightarrow$ sparsity $\rightarrow$ pruning 30\% $\rightarrow$ chưng cất $\rightarrow$ QAT\\
Độ phức tạp & 2,17 M tham số, 4,82 GFLOPs\\
Precision & INT8 (TensorRT), dự phòng FP16\\
Định dạng triển khai & TensorRT \texttt{.engine} build riêng cho từng loại Jetson\\
Phân phối & Kho model Hugging Face Hub, ghim theo commit SHA\\
Độ chính xác & \cbs{mAP@50, mAP@50:95 trên tập test}\\
Hiệu năng trên thiết bị & \cbs{FPS, latency trên Jetson Nano và Orin Nano}\\
\bottomrule
\end{tabularx}
\end{table}

\section{Kết chương}

Chương 3 đã xây dựng pipeline tối ưu YOLO26n sáu giai đoạn: tìm siêu tham số
bằng SHO, thưa hóa BatchNorm, cắt tỉa kênh 30\% theo ngưỡng toàn cục, phục hồi
bằng chưng cất CWD hoặc fine-tune, QAT với hiệu chỉnh entropy và chuyển đổi
TensorRT INT8. Cắt tỉa 30\% giảm 15,7\% tham số và 21,2\% FLOPs. Mỗi giai đoạn là
một script độc lập có thể chạy lại, đầu ra được lưu cùng cấu hình để truy vết.
Model sau tối ưu là một tệp engine được phân phối qua kho model và triển khai bởi
nền tảng được thiết kế ở Chương 4.

